% Protocol transcribed from context/08_implementation_ground_truth.md.
% Algorithms are introduced beside their explanation; floats may share pages with prose.
\section{Client-assisted CKKS protocol}
\label{sec:protocol}

Table~\ref{tab:notation} defines the evaluated configuration. Figure~\ref{fig:architecture} locates
the server computation and the client calls within a transformer block. The protocol retains the
released weights and exact nonlinearities while bounding encrypted depth through fresh client
encryption.

\begin{table}[!t]\centering\small
\caption{Notation. All values are those of the evaluated configuration.}
\label{tab:notation}
\begin{tabularx}{\columnwidth}{l l Y}
\toprule
Symbol & Value & Meaning \\
\midrule
$D$        & 768    & hidden width of the released model \\
$H$        & 12     & attention heads, of $D/H = 64$ dimensions each \\
$D_{\text{mlp}}$ & 3072 & MLP expansion width, $4D$ \\
$T$        & 103    & tokenized prompt length for the genomic-signal task \\
$B$        & 8      & tokens packed per ciphertext (token lanes) \\
$G$        & 13     & ciphertext groups, $\lceil T/B \rceil$ \\
$W$        & 1024   & feature width padded to a power of two \\
$C$        & 4      & activation copies per ciphertext, $D_{\text{mlp}}/D$ \\
$N$        & 65536  & polynomial ring degree \\
$N/2$      & 32768  & usable plaintext slots per ciphertext, $C \cdot W \cdot B$ \\
$L$        & 13     & multiplicative depth \\
$\lambda$  & 128    & classical security level, in bits \\
\bottomrule
\end{tabularx}
\end{table}


\begin{figure*}[!t]\centering
  \includegraphics[width=\textwidth]{fig_architecture.pdf}
  \caption{Client-assisted block evaluation. The compute provider evaluates encrypted linear algebra
  and the data owner evaluates exact nonlinearities at declared boundaries. The provider has public
  weights, keys, and execution metadata, but receives no plaintext query-derived activation or secret key.}
  \label{fig:architecture}
\end{figure*}

\subsection{Client boundaries and normalization}
CKKS carries affine maps, reductions, residual additions, and attention products. The inverse
square root in LayerNorm, causal softmax, and GELU use declared client calls. At each call, the
data owner decrypts values derived from its query, evaluates the function in double precision, and
encrypts the result afresh under the same key. Algorithm~\ref{alg:boundary} gives this operation.

\begin{algorithm}[H]
\small
\caption{The client boundary. The data owner is the only party holding $\mathsf{sk}$, and every
value it decrypts derives from its own query.}
\label{alg:boundary}
\begin{algorithmic}[1]
\Function{ClientBoundary}{$c$, $f$}
  \State $v \gets \Call{Decode}{\Call{Decrypt}{\mathsf{sk},c}}$
  \State $v' \gets f(v)$ \Comment{evaluated exactly, in double precision}
  \State \Return $\Call{Encrypt}{\mathsf{pk},\Call{Encode}{v'}}$ \Comment{fresh ciphertext at
         level $0$}
\EndFunction
\end{algorithmic}
\end{algorithm}

Fresh encryption returns the activation to level~$0$. The configured depth of $13$ therefore bounds
the longest encrypted run between boundaries rather than the depth accumulated across all $12$
transformer blocks. The ciphertext lineage remains at the server between these calls.

For LayerNorm, the server computes the mean and centered variance under encryption and adds
$\varepsilon$ before sending the variance term to the client. Only the inverse square root is
evaluated there; the server then applies the returned factor and public scale, as specified in
Algorithm~\ref{alg:layernorm}. No bias is omitted as an approximation: the released checkpoint
contains no bias tensors.

\begin{algorithm}[H]
\small
\caption{LayerNorm. The server computes both statistics under encryption; only the inverse square
root crosses the boundary. The released checkpoint carries no bias term.}
\label{alg:layernorm}
\begin{algorithmic}[1]
\Function{LayerNorm}{$x$, $\gamma$}
  \State $\mu \gets \Call{ReduceSum}{x}\odot(1/D)$
  \State $z \gets x - \mu$
  \State $\sigma^2 \gets \Call{ReduceSum}{z \odot z}\odot(1/D) + \varepsilon$
  \State $r \gets \Call{ClientBoundary}{\sigma^2,\ t \mapsto t^{-1/2}}$
  \State \Return $(z \odot r) \odot \gamma$
\EndFunction
\end{algorithmic}
\end{algorithm}

\subsection{Packing and dense transforms}
The slot layout in Fig.~\ref{fig:packing} assigns $32{,}768$ slots as $4$ activation copies by
$1{,}024$ padded features by $8$ token lanes. Token lane is the innermost index: whole-lane
rotations move features without mixing tokens, and smaller rotations align token lanes. The
released width of $768$ occupies the active feature region; padding supports the power-of-two
transform and attention-score staging. Rows and columns outside the active width contribute zeros.

\begin{figure}[!t]\centering
  \includegraphics[width=\columnwidth]{fig_packing.pdf}
  \caption{Slot layout and grouping for the measured prompt. The final ciphertext group carries
  seven live tokens and one padded lane; the groups produce the lower-triangular causal score tiles.}
  \label{fig:packing}
\end{figure}

The four copies express the MLP expansion from width $768$ to $3{,}072$ as parallel chunks.
The $103$ tokens occupy $13$ ciphertext groups. A $32\times32$ baby-step giant-step decomposition
covers the $1{,}024$ padded features. Algorithm~\ref{alg:matmul} combines shared baby rotations
with encoded weight diagonals, cloning each plaintext for use at the receiving ciphertext's level.

\begin{algorithm}[H]
\small
\caption{Packed matrix product, baby-step giant-step over the packed diagonals.
$N_1 = N_2 = 32$ and $N_1N_2 = W = 1024$.}
\label{alg:matmul}
\begin{algorithmic}[1]
\Function{MatMul}{$\beta$, $M$}
  \State $d \gets \Call{EncodedDiagonals}{M,\ \mathrm{level}(\beta_0)}$
         \Comment{encoded once, cloned per use}
  \State $y \gets 0$
  \For{$j \gets 0$ \textbf{to} $N_2-1$}
    \State $u \gets \sum_{i<N_1} \beta_i \odot \Call{Clone}{d_{N_1 j + i}}$
    \State $y \gets y + (j = 0\ ?\ u : \Call{Rot}{u, B N_1 j})$
  \EndFor
  \State \Return $y$
\EndFunction
\end{algorithmic}
\end{algorithm}

Eight-token packing reduces the serial-equivalent dense-product count from $1{,}236$ to $156$,
a structural reduction of $7.92\times$. Rotations, masking, encoding, and client work remain, so
this operation-count ratio does not specify a latency speedup.

\subsection{Block evaluation}
\label{sec:algorithms}
Algorithm~\ref{alg:block} follows the released block through normalization, query/key/value
projection, causal attention, output projection, residuals, and the expanded MLP. The cached
diagonals are flushed after the last use of each weight set; this changes storage lifetime without
changing the arithmetic. Baby rotations are shared among transforms that consume the same input.

\begin{algorithm}[!t]
\small
\caption{Evaluating one transformer block. $G=13$ token groups; $\odot$ is homomorphic
multiplication; every operation outside a \textsc{Client} call runs on ciphertexts.}
\label{alg:block}
\begin{algorithmic}[1]
\Require encrypted activations $c_0,\dots,c_{G-1}$; public weights $W$
\Statex \textit{Stage 1: normalization and projection}
\For{$g \gets 0$ \textbf{to} $G-1$}
  \State $n \gets \Call{LayerNorm}{c_g,\ \gamma_1}$ \Comment{Algorithm~\ref{alg:layernorm}}
  \State $\beta \gets \Call{BabyRotations}{n}$ \Comment{$N_1-1$ rotations, shared}
  \State $Q_g \gets \Call{MatMul}{\beta,W_q}$; $K_g \gets \Call{MatMul}{\beta,W_k}$
  \State $V_g \gets \Call{MatMul}{\beta,W_v}$
\EndFor
\State \Call{FlushEncodedWeights}{} \Comment{bounds host memory to one stage}
\Statex \textit{Stage 2: causal attention scores}
\For{$k \gets 0$ \textbf{to} $G-1$}
  \State $\tilde{K} \gets \{\Call{LaneShift}{K_k,\delta}\}$ for each $\delta$ required by a
         later query group
  \For{$q \gets k$ \textbf{to} $G-1$}
    \State $S \gets 0$
    \For{$\delta \in \Call{ActiveShifts}{q,k}$}
      \State $P \gets Q_q \odot \tilde{K}_\delta$
      \For{$h \gets 0$ \textbf{to} $H-1$}
        \State $S \gets S + \Call{ReduceSum}{P \odot \mathrm{mask}_h} \odot
               \mathrm{isolate}(q,k,\delta,h)/\sqrt{d_h}$
      \EndFor
    \EndFor
    \State \Call{Client.TakeScoreTile}{$S,q,k$} \Comment{decrypt only}
  \EndFor
\EndFor
\State \Call{Client.Softmax}{} \Comment{stable softmax over complete causal rows}
\Statex \textit{Stage 3: attention context}
\For{$k \gets 0$ \textbf{to} $G-1$}
  \State $\tilde{V} \gets \{\Call{LaneShift}{V_k,\delta}\}$
  \For{$q \gets k$ \textbf{to} $G-1$; $\ \delta \in \Call{ActiveShifts}{q,k}$}
    \State $A \gets \Call{Client.EmitWeightTile}{q,k,\delta}$ \Comment{re-encrypt only}
    \State $C_q \gets C_q + A \odot \tilde{V}_\delta$
  \EndFor
\EndFor
\Statex \textit{Stages 4 and 5: projection, MLP, residuals}
\For{$g \gets 0$ \textbf{to} $G-1$}
  \State $R_g \gets c_g + \Call{MatMul}{\Call{BabyRotations}{C_g},W_{\text{proj}}}$
\EndFor
\State \Call{FlushEncodedWeights}{}
\For{$g \gets 0$ \textbf{to} $G-1$}
  \State $\beta \gets \Call{BabyRotations}{\Call{LayerNorm}{R_g,\gamma_2}}$
  \State $h \gets \sum_{c<4} \Call{MatMul}{\beta,W_{\text{fc}}^{c}} \odot \mathrm{copy}_c$
  \State $a \gets \Call{Client.GELU}{h}$
  \State $c'_g \gets R_g + \sum_{c<4}
         \Call{MatMul}{\Call{BabyRotations}{a \odot \mathrm{copy}_c},W_{\text{mlp}}^{c}}$
\EndFor
\State \Return $c'_0,\dots,c'_{G-1}$
\end{algorithmic}
\end{algorithm}

The protocol transcript is fixed by the circuit and sequence length. Each block makes $857$
physical client boundary crossings, including $91$ score-tile decryptions and $727$ weight-tile
encryptions. They carry $129{,}162$ logical nonlinear instances. Softmax accounts for $128{,}544$,
or $99.5\%$, of them, concentrating the boundary work in attention. The implementation asserts
the operation schedule, boundary counts, and remaining depth at run time and aborts on deviation.

\subsection{Causal attention}
Attention visits the $91$ lower-triangular pairs of the $13$ token groups. Shifts required by
later query groups are computed once per key group and reused. The client collects all encrypted
score tiles before stable softmax normalizes complete causal rows, then returns fresh encrypted
weight tiles for context aggregation. This normalization is a global synchronization point.

Both query--key scoring and weight--value aggregation multiply two activation-derived operands.
Of the block's $1{,}506$ ciphertext--ciphertext multiplications, LayerNorm contributes $26$
centered squares and $26$ products with returned inverse square roots. Attention contributes
$727$ query--key and $727$ weight--value products, or $97\%$ of the total.

\subsection{Depth, parameters, and composition}
The retained context has ring dimension $65{,}536$, $32{,}768$ usable slots, and a $128$-bit
classical security level. It uses $50$-bit scaling, a $60$-bit first modulus, $3$ hybrid
key-switching digits, and $80$ rotation keys. Depth $13$ is the smallest demonstrated passing
depth for these parameters and packing; the block finishes with $6$ levels remaining.

Between blocks, the data owner decrypts the full hidden state, restores its token-vector layout,
and encrypts fresh level-$0$ inputs under the same context and key lineage. The preliminary
two-block execution at two tokens reached relative error $8.48\times10^{-11}$ without GPU-memory
growth in the second block. Section~\ref{sec:results} reports the complete task-length composition.

\subsection{Non-interactive comparison}
The non-interactive design replaces every nonlinearity with a polynomial approximation and
retains a single ciphertext lineage without intermediate decryption. It completes one
released-weight block. When configured for composition, however, its context and evaluation keys
exceed the tested $80$~GB GPU-memory envelope before useful composition arithmetic begins.

Repeated polynomial evaluation requires a deeper modulus budget and a bootstrapping schedule
whose associated material must remain resident. This memory limit applies to the tested library,
parameters, packing, and hardware. The client-assisted design shortens encrypted segments by
evaluating the exact functions at the client and refreshing through fresh encryption.
